\section{Discussion and Limitations}
The empirical case for BPGS is strongest on the scale-mismatch questions the method was designed to address. The rescaling study isolates numerical-scale invariance directly, the scale-stress study tests degradation under worsening imbalance, and the NYUv2 results show that the same parameterization remains competitive on a standard dense-prediction benchmark. The paper therefore supports a targeted claim about robustness-oriented uncertainty weighting rather than a general claim of dominance over all multi-task optimizers.

The study also has clear limits. First, the benchmark set is small: one dense-prediction
benchmark, two tabular real-data problems, and controlled stress tests. A second dense-prediction
benchmark remains future work. Second, most headline results use three seeds; the 10-seed
scale-stress study strengthens one comparison but does not justify fine-grained ranking claims.
Third, BPGS trails Kendall and PCGrad in the noisy and conflict mixed-stress regimes. It
stabilizes scalar weights but does not alter gradient directions, so conflict-aware methods can
have a structural advantage when gradient interference is the dominant difficulty. Combining BPGS
with such methods is untested future work.

Fourth, boundedness does not guarantee learnability at the chart boundary. The logged runs remain
inside the boundary (the largest observed $|z_i|/\tau_T$ is 0.93), but this is an empirical margin,
not a guarantee for new tasks or training regimes. Fifth, the batch-size and first-batch studies
cover only batch sizes 4--16 and the tested NYUv2 subset. Sixth, the formal results establish
batch-conditional boundedness and normalized-weight invariance, not convergence or
scale-independent optimality. Finally, Nash-MTL is the only added recent baseline; CAGrad,
IMTL-G, FAMO, and Auto-Lambda \cite{liu2022autolambda} remain outside the evaluated comparison
set. The runtime study quantifies cost only on the controlled NYUv2 subset and one GPU setting.

Potential societal impact is indirect. A more stable task-weighting rule can improve the reliability of multi-output models used in scientific or engineering pipelines, but any downstream deployment still inherits the data, labeling, and evaluation failures of the underlying tasks. In safety-sensitive settings, a method that improves one task while weakening another can be misused if only aggregate performance is reported, which is why this paper reports per-task metrics and states where BPGS is not the best method.
